% \documentclass{report}

% \usepackage{graphicx}
% \usepackage{dcolumn}
% \usepackage{bm}
% \usepackage{amsfonts, amssymb, amsmath}
% \usepackage{braket}

%\newcommand{\ket}[1]{\textstyle \left| #1 \right\rangle \displaystyle}
%\newcommand{\bra}[1]{\textstyle \left\langle #1 \right| \displaystyle}
%\newcommand{\braket}[2]{\textstyle \left\langle #1 \left|  #2 \right\rangle\right. \displaystyle}
%\newcommand{\im}{\mathrm{i}}
%\newcommand{\e}{\mathrm{e}}
%\newcommand{\pwrt}[2]{\frac{\partial #1}{\partial #2}}
%\newcommand{\diff}{\mathrm{d}}

%\begin{document}

\subsection{Symmetry concerns within silicon}
\label{subsection:symm_silicon}

The operations of the point-group symmetry of our system, a silicon donor, can
be conveniently grouped into a few categories. These categories are defined by
their transformation effects upon functions with crystalline Cartesian
coordinate domain. For convenience, we'll denote the function transformation
operators compactly such that $\left(\bar{y}x\bar{z}\right)$ indicates an
operator acting such that $\left(\bar{y}x\bar{z}\right) f_{xyz} =
f_{\bar{y}x\bar{z}}$. For our purposes we recognize 
\begin{itemize}
  \item cyclic permutation e.g. $\left(zxy\right)$,
  \item acyclic swap e.g. $\left(yxz\right)$,
  \item double inversion e.g. $\left(x\bar{y}\bar{z}\right)$.
\end{itemize}
An operation in the point group can be part of two categories. An operator in
two categories will be a member of the double inversion category, as well as
either the cyclic or acyclic category. An example of such is
$\left(\bar{y}x\bar{z}\right)$. These categories form the basis of our
particular representation of the symmetries of silicon. Within our formulation
of the group, most of the representations of the $T_{d}$ point symmetry group of
silicon can be considered as having definite even or odd symmetry under acyclic
swaps and double inversion. For example, functions belonging to the $A_{1},
A_{2}$ representations are invariant under double inversion, but are
respectively classified as even or odd under acyclic swap. In comparison,
functions belonging to the $T_{1}, T_{2}$ representations are odd under double
inversion, but are respectively classified as odd or even under acyclic swap.
Note that $A_{1}, T_{2}$ are even and $A_{2}, T_{1}$ are odd under acyclic swap. The final representation, $E$, is invariant under double inversion, but cannot be said to be simply even or odd under acyclic swaps. 

This categorization of the operators of the point-group of our system allows a further useful observation; namely, we may factor the projection operator based upon these categories. To wit,
\begin{equation}
  \mathcal{P}_{ii}^{\left(\kappa\right)}
  =
  \mathcal{V}_{\kappa,i}
  \mathcal{A}_{z}^{\left(\kappa\right)}
  \mathcal{D}_{z}^{\left(\kappa\right)}
  \mathcal{V}_{\kappa,i}^{\dagger}
  ,
\label{symm-operator}
\end{equation}
where the factored operators have been determined to be
\begin{align}
  \mathcal{V}_{\kappa,i}
  & \equiv
  \alpha_{0}^{\left(\kappa,i\right)}\left(xyz\right)
  + \alpha_{1}^{\left(\kappa,i\right)}\left(yzx\right)
  + \alpha_{2}^{\left(\kappa,i\right)}\left(zxy\right),
  \label{indexing_proj}
  \\
  \mathcal{A}_{z}^{\left(\kappa\right)}
  & \equiv
  \frac{1+\delta_{E\kappa}}{2}
  \left\{
    \left(xyz\right) + \sigma^{\left(\kappa\right)}\left(yxz\right)
  \right\}
  \label{acyclic_proj}
  \\
  \mathcal{D}_{z}^{\left(\kappa\right)}
  & \equiv
  \frac{1}{4}
  \left\{
    \left(xyz\right)
    +
    \left(\bar{x}\bar{y}z\right)
    +
    \omega^{\left(\kappa\right)}
    \left(\bar{x}y\bar{z}\right)
    +
    \omega^{\left(\kappa\right)}
    \left(x\bar{y}\bar{z}\right)
  \right\}
  .
  \label{dblinv_proj}
\end{align}

Above, $\kappa$ denotes the symmetry representation, while $i$ indexes
degenerate states. The $A_{1}$ and $A_{2}$ states are nondegenerate, and do not
require indexing. The $E$ states are doubly degenerate and the $T_{1}$ and
$T_{2}$ states are triply degenerate. With our particular irreducible
representation for the $T$ symmetries, each of the triply degenerate $T$ states
can be indexed by the three valley axes. The values of the $\omega$ and $\sigma$
constants in the expressions above follow from the discussion below. 
% \begin{equation} \omega_{qr}^{\left(\kappa,i\right)}
%   =
%   \omega_{q}^{\left(\kappa,i\right)} \omega_{r}^{\left(\kappa,i\right)}
%   \end{equation}
Namely for the states where $\kappa = A_{1},A_{2},E$ we have $ \omega^{
\left(\kappa\right) } = +1 $ and $\kappa = T_{1},T_{2}$ we have $ \omega^{
\left(\kappa\right) } = -1 $. 
% For $T$ states the value of $\omega_{q}$ varies between the three degenerate $T$
% states as
% \begin{equation}
%   \omega_{q}^{ \left(T_{n}, i\right) }
%   = -\delta_{iq}+\delta_{ir}+\delta_{is}.
% \end{equation}
For $A_{1}, T_{2}$ states we find $\sigma = +1$ and for $A_{2}, T_{1}$ states we
find $\sigma = -1$. The $E$ states have $\sigma = 0$. Finally the valley axis
combination vector given by $\bm{\alpha}$ is determined for $A$, $E$, and $T$
states respectively. We find,
\begin{equation}
  \bm{\alpha}^{\left(A_{n}\right)}
  = \frac{1}{3}\begin{pmatrix} 1 \\ 1 \\ 1 \end{pmatrix},
\end{equation}
\begin{equation}
  \bm{\alpha}^{\left(E,0\right)}
  = \frac{1}{3}\begin{pmatrix} 1 \\ \tau^{2} \\ \tau \end{pmatrix},
\quad\quad
  \bm{\alpha}^{\left(E,1\right)}
  = \frac{1}{3}\begin{pmatrix} 1 \\ \tau \\ \tau^{2} \end{pmatrix},
\end{equation}
and finally,
\begin{equation}
  \bm{\alpha}^{\left(T_{n},x\right)} = \begin{pmatrix} 0 \\ 1 \\ 0\end{pmatrix},
\quad\quad
  \bm{\alpha}^{\left(T_{n},y\right)} = \begin{pmatrix} 0 \\ 0 \\ 1\end{pmatrix},
\quad\quad
  \bm{\alpha}^{\left(T_{n},z\right)} = \begin{pmatrix} 1 \\ 0 \\ 0\end{pmatrix}.
\end{equation}
Above $\tau = \mathrm{e}^{\im\frac{2\pi}{3}}$, one of the third roots of unity.
We will explore further how these symmmetry effects from the silicon media
influence the development of our donor electron's basis state.